\section{Transition-system specifications with negative premises}
\label{sec:p09-negative-premises}

A transition-system specification (TSS) defines the transition relation,
rather than taking it as input to a bisimulation or modal analysis.
The present construction follows the separate 133-page author presentation
of Aceto, Fokkink and Verhoef, \S\S2.5--3.5
\parencite{AcetoFokkinkVerhoefSOSAuthorNP}, and the ordinal construction
of \textcite[\S2]{Groote1993NegativePremises}. The author presentation
is not identified with the unread publisher chapter, pp.197--292 of
the 2001 \emph{Handbook of Process Algebra}. Predicates, arbitrary
premise sets and the distinction between supported and stable meanings
are retained. A negative premise is a statement about absence in the
specified relation, not failure of an unspecified search procedure.

\begin{hbdefinition}{def:p09-np-language}{Terms, judgements and denial}
Let \(\Sigma\) be a set-sized ranked signature, \(V\) an infinite
variable supply, \(C=T(\Sigma)\) its finite closed terms, \(A\)
a set of action labels, and \(\mathsf{Pred}\) a set of unary
predicate names. Positive closed judgements are \(t\xrightarrow a u\)
and \(tP\). Their disjoint union is the set \(B\).
Standard negative premises are \(t\mathrel{\not\!\xrightarrow a}\)
and \(t\neg P\). For a negative literal \(\eta\), let
\(D(\eta)\subseteq B\) be its denying positive judgements:
\[
 D(t\mathrel{\not\!\xrightarrow a})=
       \{t\xrightarrow a u:u\in C\},\qquad
 D(t\neg P)=\{tP\}.
\]
An interpretation \(M\subseteq B\) satisfies a positive literal
exactly when it belongs to \(M\), and satisfies a negative literal
exactly when \(M\cap D(\eta)=\varnothing\).
Ground substitution instances replace every rule variable by a closed
term. No rewrite-rule variable-containment condition is imposed on
general TSS rules. A rule has a set of positive and negative premises
\(H\) and a positive conclusion \(\alpha\), written \(H/\alpha\).
A TSS \(T\) is a set of such rules; premise sets may be infinite.

For the proof comparison only, also allow the target-specific negative
\(t\mathrel{\not\!\xrightarrow a}u\), with denying set
\(\{t\xrightarrow a u\}\). This is the author's explicit
\S3.2 proof-language extension, not an unannounced change to the
standard TSS premise language. Two literals deny each other when one
is negative and the other belongs to its denying set.
\end{hbdefinition}

\begin{hbdefinition}{def:p09-np-models}{Model, support, positive proofs and reduct}
\(M\) is a model if every ground rule with all premises satisfied
in \(M\) has its conclusion in \(M\). It is supported if every
member has such a ground-rule witness. A supported model has both
properties. Support alone permits circular positive justification.
A proof of \(H/\alpha\) is a set-sized upwardly branching tree
with no infinite branch: its root is \(\alpha\), and each node
is either a hypothesis leaf from \(H\), or a ground rule application.
Write \(T\vdash H/\alpha\). Branching need not be finite.

For a candidate \(M\), retain precisely the ground rules whose
negative premises hold in \(M\), and erase those negative premises.
This positive reduct is \(T^M\). Let \(\Gamma_T(M)\) be its
least model. A stable model satisfies \(M=\Gamma_T(M)\).
A well-supported model is a model such that every \(\alpha\in M\)
has \(T\vdash N/\alpha\) for a set of negative hypotheses \(N\)
satisfied in \(M\). These are properties of relations on the same
ground judgement set, not alternative notions of a term or an action.
\end{hbdefinition}

\begin{hblemma}{thm:p09-np-positive}{proved=annex:p09-np-positive}{Positive closure with arbitrary premises}
A positive TSS has a least model, exactly its well-founded proof-tree
conclusions. It is supported. If \(\lambda\) is an infinite regular
cardinal strictly larger than every positive premise-set cardinality,
iteration from the empty relation reaches this least model by the
union of stages below \(\lambda\). Finite premise sets permit
\(\lambda=\omega\); arbitrary premise sets do not.
\end{hblemma}

\begin{hblemma}{thm:p09-np-stable}{proved=annex:p09-np-stable}{Stable and well-supported models coincide, not merely supported models}
For every \(M\subseteq B\), \(\Gamma_T(M)\) consists exactly
of the judgements admitting a proof \(T\vdash N/\alpha\) with
negative hypotheses true in \(M\). Stable models are exactly
well-supported models. Every stable model is a supported model,
but the converse fails, even for a positive TSS.
\end{hblemma}

\begin{hbdefinition}{def:p09-np-three-valued}{The well-founded three-valued interpretation}
Use a pair \(L\subseteq U\subseteq B\): \(L\) is certainly
true, \(U\setminus L\) unknown, and \(B\setminus U\) false.
It is three-valued stable when
\(L=\Gamma_T(U)\) and \(U=\Gamma_T(L)\).
This interval convention is equivalent to the source's pair of
disjoint true/unknown sets \(\langle C,U_{\rm unknown}\rangle\).
The well-founded interpretation is
\[
 L_* = \mu X\,\Gamma_T(\Gamma_T(X)),\qquad
 U_* = \Gamma_T(L_*).
\]
The least fixed point is a powerset least fixed point, with transfinite
iteration if needed. It is not a finite failure test.
\end{hbdefinition}

\begin{hbtheorem}{thm:p09-np-well-founded}{proved=annex:p09-np-well-founded}{Existence and the exact information order}
\(\Gamma_T\) is antitone; its square is monotone. The displayed
well-founded pair exists, satisfies \(L_*\subseteq U_*\), and
is the information-least three-valued stable pair: for any other
\(L\subseteq U\), \(L_*\subseteq L\) and \(U\subseteq U_*\).
Every stable model lies between its two bounds. A total well-founded
pair gives a unique stable model, but unique stability alone need
not make the well-founded pair total.
\end{hbtheorem}

\begin{hbdefinition}{def:p09-np-proof-meanings}{Supported and well-supported literal proofs}
Both literal-proof notions use well-founded trees. Positive nodes use
ground TSS rules. At a negative node \(\eta\), let \(K\) be
its children's labels. A supported proof must block every ground
rule with conclusion denying \(\eta\), by finding a premise of
that rule denying a member of \(K\). A well-supported proof must
instead block every ordinary proof \(T\vdash N/\gamma\), where
\(N\) is negative and \(\gamma\) denies \(\eta\), by finding
a member of \(N\) denying a member of \(K\).
Their judgements are \(T\vdash_s\ell\) and
\(T\vdash_{ws}\ell\). A TSS is \(ws\)-complete when every
positive atom is either \(ws\)-provable or its target-specific
or predicate negative is \(ws\)-provable. The registered sheet
makes these universal blocking conditions explicit.
\end{hbdefinition}

\begin{hbtheorem}{thm:p09-np-proof-meanings}{proved=annex:p09-np-proof-meanings}{Proof soundness, consistency and completeness meaning}
Supported proofs are sound in every supported model.
Well-supported proofs are sound in every three-valued stable pair:
positive conclusions belong to its lower bound, and a negative
conclusion has no denier in its upper bound. Consequently no TSS
has denying \(ws\)-provable literals. In the well-founded pair,
\[
 T\vdash_{ws}\alpha\iff\alpha\in L_*,\qquad
 T\vdash_{ws}\eta\iff D(\eta)\cap U_*=\varnothing.
\]
Thus \(ws\)-completeness is exactly totality of the well-founded
pair. Merely ruling out multiple stable models is a different criterion.
\end{hbtheorem}

\begin{hbdefinition}{def:p09-np-stratification}{Ordinal stratification over every possible ground judgement}
A stratification is \(S:B\to\kappa\), for an ordinal \(\kappa\),
such that every ground rule with conclusion \(\alpha\) satisfies
\(S(\beta)\leq S(\alpha)\) for each positive premise \(\beta\),
and \(S(\gamma)<S(\alpha)\) for \emph{every}
\(\gamma\in D(\eta)\) of each negative premise \(\eta\).
In particular an action-absence premise quantifies over all closed
targets, not only transitions already derived. A strict stratification
also makes each positive-premise inequality strict.

Let \(B_\beta=S^{-1}(\beta)\),
\(I_{<\beta}=\bigcup_{\gamma<\beta}I_\gamma\).
After the lower strata are fixed, define the monotone operator
\[
 F_\beta(X)=\{\alpha\in B_\beta:\text{some ground rule for }\alpha
 \text{ has negatives true in }I_{<\beta}
 \text{ and positives in }I_{<\beta}\cup X\}.
\]
Take \(X_{\beta,0}=\varnothing\),
\(X_{\beta,\xi+1}=F_\beta(X_{\beta,\xi})\), and unions at limits.
For one infinite regular \(\lambda\) bounding every positive
premise-set size strictly, put
\(I_\beta=\bigcup_{\xi<\lambda}X_{\beta,\xi}\),
and \(M_S=\bigcup_{\beta<\kappa}I_\beta\).
All signatures, rules and judgement universes are sets; the proof
uses ordinary classical set theory with Choice. No finiteness,
countability, image-finiteness or effective stratification is assumed.
\end{hbdefinition}

\begin{hbtheorem}{thm:p09-np-stratified}{proved=annex:p09-np-stratified}{The full stratified supported relation and its stronger semantic status}
The construction is well-defined and gives a supported model,
the unique stable model, and the total well-founded interpretation.
It is \(ws\)-complete. Its result is independent of the valid
stratification and regular bound chosen. It need not be the unique
supported model; under \emph{strict} stratification it is.
The same-stage positive closure and the regular-cardinal bound
are essential parts of the argument, not unnamed dependencies.
\end{hbtheorem}

\begin{hbdefinition}{def:p09-np-countermodels}{The source's six small TSSs and two further boundaries}
For the predicate-only examples put \(h=aP_1\), \(k=aP_2\),
with one source constant \(a\), the shown predicates and no
action judgements. A bar denotes the corresponding predicate negative.
The registered \(T_1\)--\(T_6\) rules are
the author's actual examples on pp.18--21. Their interpretation
bases contain only \(h\), or \(h,k\), as appropriate.
Also register an editorial three-predicate example \(E\), with
an additional \(r=aP_3\), and an infinitary positive transition
example \(J\), with one state \(f\), labels \(n\in\mathbb N\)
and a distinct label \(b\). These latter examples are not attributed
as printed source examples.
\end{hbdefinition}

\begin{hbtheorem}{thm:p09-np-countermodels}{proved=annex:p09-np-countermodels}{Exact absence, ambiguity, circular-support and ordinal countercases}
\(T_1\) has two stable models \(\{h\},\{k\}\) and no
least model. \(T_2\) has least model \(\{h\}\) but no supported
or stable model. \(T_3\) has least supported/stable model
\(\{h\}\), but no least model. \(T_4\) has least model
\(\{h\}\), unique supported model \(\{h,k\}\), and no
stable model. \(T_5\) has supported models \(\varnothing,\{h\}\)
but unique stable/well-founded model \(\varnothing\).
\(T_6\) has supported models \(\{h\},\{k\}\)
and unique stable/well-founded model \(\{k\}\).
For \(E\) the unique stable model is \(\{h\}\), but all three
judgements are unknown in the well-founded interpretation.
In \(J\), an infinite positive premise set requires a judgement
appearing beyond stage \(\omega\); truncation to finite stages
does not even give a model. The annex supplies the actual calculations.
\end{hbtheorem}

\subsection{The source process rules and actual stratifications}

The sheet retains Tables 1--3 of the author source: BPA with empty
process \(\epsilon\), alternative and sequential composition;
the unary priority operator \(\theta\); and the separate relative
discrete-time extension with delay \(\sigma_d\) and time label
\(\sigma\notin A\). Successful termination is the predicate
\(t\checkmark\), not absence of transitions. The priority guard
contains every \(b\) with \(a<b\) for the source's strict partial
order on actions. The time-choice rules have the source's separate
both-can-delay and exactly-one-can-delay branches.

\showcalculus{p09-negative-premises}

\begin{hbtheorem}{thm:p09-np-processes}{proved=annex:p09-np-processes}{All source process rules, ranks and small transition calculations}
Each of the three separate source TSSs has a total well-founded and
unique stable relation. For priority, count \(\theta\)-occurrences
in the source of every judgement. For discrete time, rank timed
judgements by the number of choice occurrences in their source,
and ordinary action and termination judgements by 0. These are the
actual source stratifications; every positive and negative condition
is checked in the annex. Source-size also supplies a strict
stratification for these nonrecursive term languages, hence their
supported models are unique as well.

For \(A=\{a,b\}\) with \(a<b\), \(\theta(a+b)\) has
exactly the \(b\)-transition to \(\theta(\epsilon)\), and no
\(a\)-transition. In the separate timed TSS,
\(\sigma_d(a)+b\) has a time transition exactly to \(a\),
whereas \(\sigma_d(a)+\sigma_d(b)\) has a time transition
exactly to \(a+b\). The one-branch rules do not apply to the
second example. These are exact syntactic transitions, not claims
of timed-automaton region equivalence.
\end{hbtheorem}

\begin{proofboundary}{def:p09-np-boundaries}{Construction is not a congruence theorem or a whole-source substitute}
The resulting action relation is an LTS on closed terms; predicates
give proposition valuations for the existing modal interface.
Neither totality of the action relation nor finiteness of successor
sets follows. Bisimulation congruence requires a separate rule-format
theorem; stratification alone does not supply one. Recursive priority,
abstraction, conservative-extension formats, higher-order/many-sorted
rules and the remaining publisher/author source components retain
their separate hypotheses and proofs. No source unit is completed
merely by this contribution to planned section 133.2.
\end{proofboundary}
